\documentclass[a4paper,11pt]{article}
\usepackage[margin=24mm]{geometry}
\usepackage[T1]{fontenc}
\usepackage{lmodern}
\usepackage{amsmath,amssymb,amsthm,mathtools}
\usepackage{microtype,booktabs}
\usepackage[colorlinks=true,linkcolor=blue,urlcolor=blue]{hyperref}
\usepackage{bookmark}
\emergencystretch=2em
\allowdisplaybreaks[1]
\newtheorem{theorem}{Theorem}[section]
\newtheorem{lemma}[theorem]{Lemma}
\newtheorem{proposition}[theorem]{Proposition}
\newtheorem{corollary}[theorem]{Corollary}
\theoremstyle{remark}
\newtheorem{remark}[theorem]{Remark}
\newcommand{\E}{\mathbb E}
\newcommand{\R}{\mathbb R}
\newcommand{\Pp}{\mathbb P}
\newcommand{\norm}[1]{\lVert#1\rVert}
\newcommand{\Ltwo}[1]{\lVert#1\rVert_{L^2}}
\title{Finite Comparison of Initialized Private Adam\\under Batch Scaling}
\author{}
\date{16 November 2026}
\begin{document}
\maketitle
\begin{abstract}
We give an explicit finite comparison of expected squared-gradient risks for
two initialized private Adam runs with different Poisson batch sizes. The
learning rate is multiplied by the square root of the batch multiplier,
the momentum coefficients use the linear transformation
$\beta_i'=1-k(1-\beta_i)$, and the outside-root stabilizer is divided by
$\sqrt{k}$. Both finite bias corrections and every original output age are
retained. A direct coupling to a common stochastic differential equation gives
an error of order $\sqrt h(1+\log(1/h))$ for fixed sampling rates and fixed scaling
parameters. The proof controls the random denominator jointly with its
numerator and accounts for the finite sampling fluctuations. Each run has a
separate differential-privacy guarantee under an explicit public calibration.
A finite example gives risk difference below $10^{-6}$, at an extremely large
horizon. The conclusion is closeness, not an ordering or an optimal choice of
learning rate. Uniformity at bounded epoch budgets is not asserted.
An additional expected-memory estimate improves the linear-to-exponential
memory-transfer contribution from order $\sqrt h$ to order $h$; it does not
change the order of the complete comparison.
\end{abstract}

\section{Mechanisms and the finite batch comparison}
\label{sec:main}

Fix a public finite universe of records. Each record has a globally Lipschitz
gradient $g_r:\R^p\to\R^p$ with $\norm{g_r(x)}\le1$ for all $x$.
Datasets are subsets, with add/remove adjacency. Fix a positive public
integer $N$ independently of the private input. The utility statement is
evaluated on a reference dataset of size $N$, with objective gradient
$\nabla f=N^{-1}\sum_{r=1}^N g_r$. Write $L_r$ for the record Lipschitz
constants and take $L>0$ with $N^{-1}\sum_rL_r\le L$. In particular
$\norm{\nabla f}\le1$ and $\nabla f$ is $L$-Lipschitz.
The initial position $x_0$ and all algorithm parameters are public.

At each age $j$ use independent Bernoulli$(q)$ record indicators $I_{j,r}$,
with the same indicator in every coordinate of a record, and a fresh standard
$p$-dimensional Gaussian $Z_j$, independent of all masks. For a public
normalizer $B=qN$, the release and Adam recursion are
\begin{align}
 Y_j&=\frac1B\sum_{r\in D}I_{j,r}g_r(x_j)+\frac\sigma B Z_j,
       \label{eq:release}\\
 m_{j+1}&=\beta_1m_j+(1-\beta_1)Y_j,&
 v_{j+1}&=\beta_2v_j+(1-\beta_2)Y_j^{\odot2},\notag\\
 \widehat m_{j+1}&=\frac{m_{j+1}}{1-\beta_1^{j+1}},&
 \widehat v_{j+1}&=\frac{v_{j+1}}{1-\beta_2^{j+1}},\notag\\
 x_{j+1}&=x_j-\eta\frac{\widehat m_{j+1}}
                         {\sqrt{\widehat v_{j+1}}+\zeta},&m_0&=v_0=0.
       \label{eq:adam}
\end{align}
All vector norms are Euclidean. Roots, squares and division are coordinatewise. Unit clipping is inactive
under the record bound. The normalizer $B$ is not recomputed on a neighboring
dataset. The masks remain internal randomness.
The privacy statement quantifies over all neighboring subsets of the public
universe using this same fixed $N$ and the same public parameters; it does
not disclose or recompute the input cardinality.

For a run of $T$ updates, its risk is
\begin{equation}\label{eq:risk}
 R=\frac1T\sum_{j=0}^{T-1}\E\norm{\nabla f(x_j)}^2.
\end{equation}
This is the risk at an independent uniform original output age. Releasing
that parameter is postprocessing; releasing an additional dataset-dependent
gradient evaluation is not asserted to be postprocessing.

\paragraph{Two exact clocks.}
Fix $H,e_0,\lambda_1,\lambda_2>0$ with $\lambda_2<2\lambda_1$,
integers $k,n\ge1$, and $0<q\le1/k$. Put $h=H/(kn)$ and
\begin{equation}\label{eq:pair}
\begin{array}{c|cc}
 &\text{source}&\text{target}\\ \hline
 T&kn&n\\
 q_i&q&kq\\
 B_i&qN&kqN\\
 \eta_i&\sqrt h&\sqrt{kh}\\
 \zeta_i&e_0/\sqrt h&e_0/\sqrt{kh}\\
 \beta_{r,i}&e^{-\lambda_rh}&1-k(1-e^{-\lambda_rh})
\end{array}
\end{equation}
The expected numbers of passes, $T_iq_i$, agree. The target momentum in
\eqref{eq:pair} is linear in the source momentum, rather than its $k$th power.

Fix $0<\epsilon\le1$, $0<\delta<1$, and the public order
\begin{equation}\label{eq:privacy-cal}
 a_{\rm priv}=\left\lceil\frac{2\log(1/\delta)}\epsilon\right\rceil+1,
 \qquad \sigma_i^2=\frac{2T_i a_{\rm priv}q_i^2}{\epsilon}.
\end{equation}
Thus $\sigma_t=\sqrt{k}\sigma_s$. Define
\begin{gather}
 c=\frac{\sqrt{2Ha_{\rm priv}/\epsilon}}N,\qquad
 K=e^{|\lambda_1-\lambda_2|H/2},\qquad
 C_\beta=H(k-1)(\lambda_1^2+\lambda_2^2/2),\notag\\
 L_q=\frac{\sum_rL_r}{kqN},\qquad
 C_x=\frac{(1+K)L_q}{e_0},\qquad
 A_{\rm mem}=H\sqrt p K C_\beta e^{C_\beta\bar h+HC_x}.
 \label{eq:memory-constants}
\end{gather}
Here $\bar h>0$ is a fixed public upper bound for $h$.
Section~\ref{sec:coupling} defines a completely explicit function
$\mathcal E(\tau;M,L_0,L_\phi)$, with the other parameters
$H,p,\lambda_1,\lambda_2,c,e_0$ fixed.

\begin{theorem}[Finite private batch comparison]\label{thm:pair}
Assume the definitions above and
\begin{equation}\label{eq:guards}
 0<h\le\bar h\le\frac1{2k\max(\lambda_1,\lambda_2)},\qquad
 h\le\min\left\{
 \frac{2\lambda_1-\lambda_2}{2(k+1)\lambda_1^2},
 \frac{4Hq^2}{\epsilon(a_{\rm priv}-1)}\right\}.
\end{equation}
The two runs in \eqref{eq:pair} are well-defined and each transcript and
independent uniform-age parameter output is separately
$(\epsilon,\delta)$-differentially private. Their actual risks satisfy
\begin{equation}\label{eq:main-error}
 |R_t-R_s|\le
 2L A_{\rm mem}\sqrt{h/k}
 +\mathcal E(h;1/q,L,2L)
 +\mathcal E(kh;1/(kq),L,2L).
\end{equation}
For fixed $H,p,N,q,k,e_0,\lambda_i,\epsilon,\delta$ and gradient bounds,
the right side is $O(\sqrt h(1+\log(1/h)))$ as $n\to\infty$.
\end{theorem}

Appendix~\ref{app:expected-memory} further permits the smaller of the
pathwise memory contribution in \eqref{eq:main-error} and an explicit
expected-memory contribution of order $h$.

The conclusion compares expected risks, rather than sufficient upper bounds
on those risks. It does not determine their order or minimize either risk.
Each finite pair has matching epochs, but the common epoch count grows with
$n$. A numerical finite example appears in Section~\ref{sec:example}.

\paragraph{Structure of the argument.}
For exponential moment parameters, the second memory after multiplication
by the mesh concentrates about $c^2$. Its fluctuation is used jointly with
the first moment through Cauchy--Schwarz. Replacing the denominator in an
auxiliary comparison, and paying the centered sampling sum, leaves a recursion
with a deterministic denominator. Its corrected first moment is an exponential
convolution. Coupling its Gaussian increments to Brownian increments gives
an explicit comparison with an integrable, singular-at-zero convolution.
Finally, normalized memory weights compare the linear target to the
exponential target. The latter and the source have the same limiting law.

\section{An explicit initialized coupling}\label{sec:coupling}

Consider a single exponential run with mesh $\tau=H/n_\tau$,
$\beta_i=e^{-\lambda_i\tau}$, $\eta=\sqrt\tau$,
$\zeta=e_0/\sqrt\tau$, and release
$Y_j=A_j(x_j)+cZ_j/\sqrt\tau$. Let $\mathcal F_j$ denote the history
before the current sampling. More explicitly, take sigma fields
$\mathcal F_j\subseteq\mathcal G_j\subseteq\mathcal F_{j+1}$ containing
all earlier signals and Gaussian innovations, with $x_j$ measurable for
$\mathcal F_j$, the sampled signal $A_j(x_j)$ measurable for
$\mathcal G_j$, and $Z_j$ standard Gaussian independent of
$\mathcal G_j$. All evaluated random variables are measurable. In the
Poisson construction, $\mathcal G_j$ additionally reveals the current
record mask. Assume globally
\begin{equation}\label{eq:oracle-assumptions}
 \norm{A_j(x)}\le M,\quad \norm{b(x)}\le M,\quad
 \E[A_j(x_j)\mid\mathcal F_j]=b(x_j),\quad
 \norm{b(x)-b(y)}\le L_0\norm{x-y}.
\end{equation}
The current signal is revealed before the independent Gaussian.
The separate global bound on $b$ is needed because comparison paths can
visit positions other than the sampled trajectory. For Poisson sampling
on the reference dataset, $b=\nabla f$, $M=1/q$ and $L_0=L$ suffice.

Write $D=c+e_0$, $E_i=e^{\lambda_iH}$, $a_i=\sqrt{E_i}$ and define
\begin{align}
 B_\tau&=a_2(\sqrt2c^2+2cM\sqrt\tau),\notag\\
 F_\tau&=\frac p{ce_0D}\left[
 2MB_\tau\sqrt{H\tau}
 +ca_1B_\tau\sqrt\tau(1+\log n_\tau)
 +HM^3\tau+2ca_1M^2\tau\sqrt H\right],\notag\\
 S_\tau&=\frac{2Ma_1\sqrt{H\tau}}D,\qquad
 A_\tau=e^{L_0H/D}(F_\tau+S_\tau),\notag\\
 Q_\tau&=2M\tau(1+E_1\log n_\tau)
 +\sqrt p c\sqrt\tau\{E_1(3+\log n_\tau)+2E_1^2\}
 +2\sqrt p c\lambda_1E_1\tau\sqrt H,\notag\\
 \omega_\tau&=\frac{M\tau+2\sqrt p cE_1\sqrt\tau}D,
 \qquad C_\tau=e^{L_0H/D}\left(\frac{Q_\tau}D+
                              \frac{L_0H\omega_\tau}D\right),\notag\\
 \mathcal E(\tau;M,L_0,L_\phi)&=L_\phi(A_\tau+C_\tau+\omega_\tau).
 \label{eq:constants}
\end{align}
These formulas also apply when $n_\tau=1$, with $\log n_\tau=0$.

Let $B_t$ be standard $p$-dimensional Brownian motion. Define the continuous
process by
\begin{align}
 m(t)&=\lambda_1\int_0^t e^{-\lambda_1(t-s)}b(X_s)\,ds
       +\lambda_1c\int_0^t e^{-\lambda_1(t-s)}\,dB_s,\notag\\
 X_t&=x_0-\frac1D\int_0^t\frac{m(s)}{1-e^{-\lambda_1s}}\,ds.
 \label{eq:limit}
\end{align}
The integrand's value at zero is immaterial.
We use a completed filtration and a jointly measurable version of the
Gaussian convolution whose unnormalized paths are continuous. All assertions
about its absolute time integrability are almost sure assertions. The
resulting continuous position process starts at the stated deterministic
$x_0$.

\begin{theorem}[Single-run quantitative comparison]\label{thm:coupling}
Under \eqref{eq:oracle-assumptions}, \eqref{eq:limit} has a unique continuous
adapted solution. There is a coupling of this solution with the initialized
Adam run such that
\begin{equation}\label{eq:grid-error}
 \E\max_{0\le j\le n_\tau}\norm{x_j-X_{j\tau}}
 \le A_\tau+C_\tau.
\end{equation}
For every $[0,1]$-valued $L_\phi$-Lipschitz function $\phi$,
\begin{equation}\label{eq:observable-error}
 \left|\frac1{n_\tau}\sum_{j<n_\tau}\E\phi(x_j)
          -\frac1H\int_0^H\E\phi(X_t)\,dt\right|
 \le\mathcal E(\tau;M,L_0,L_\phi).
\end{equation}
\end{theorem}

\begin{proof}
Appendix~\ref{app:moments} compares Adam to the mean recursion with constant
denominator, at cost $A_\tau$. Appendix~\ref{app:sde} constructs the solution
and compares that recursion to its grid values, at cost $C_\tau$.
The same appendix proves
$\E\norm{X_t-X_s}\le\omega_\tau$ whenever $0\le t-s\le\tau$.
Apply the Lipschitz bound first at the coupled grid values, and then between
$j\tau$ and $t\in[j\tau,(j+1)\tau]$, and integrate. This includes age zero
and does not introduce $x_{n_\tau}$ into the risk average.
\end{proof}

\begin{proof}[Proof of Theorem~\ref{thm:pair}]
Introduce a target exponential run, replacing only its momentum parameters
by $e^{-\lambda_i kh}$. Its averaged noise is $c/\sqrt{kh}$, while the
source averaged noise is $c/\sqrt h$, by \eqref{eq:privacy-cal}.
Both mean gradients are $\nabla f$. They therefore have the same uniquely
defined law \eqref{eq:limit}. One may realize this comparison directly on
one Brownian motion and independent source and target mask arrays: use
$h$-increments for the source, $kh$-increments for the target, and the same
target masks and increments for its linear and exponential versions.
Each run has its prescribed marginal law. Both exponential versions then
compare to the same pathwise solution $X$. This mathematical coupling makes
no assertion about privacy of a joint release.
Apply Theorem~\ref{thm:coupling} to both,
using $M_s=1/q$, $M_t=1/(kq)$ and
$\phi(x)=\norm{\nabla f(x)}^2$. This function is $2L$-Lipschitz and
bounded by one. Appendix~\ref{app:memory} compares the linear and exponential
target positions at every age with error $A_{\rm mem}\sqrt{h/k}$.
The triangle inequality through the exponential target and common limit
proves \eqref{eq:main-error}.

For positivity and separation, put $d_i=1-e^{-\lambda_i h}$ and
$\Delta=2\lambda_1-\lambda_2>0$. The bounds
$u-u^2/2\le1-e^{-u}\le u$ give
\[
 1-kd_i\ge\tfrac12,\qquad
 (1-kd_2)-(1-kd_1)^2
 =k(2d_1-d_2-kd_1^2)
 \ge k\{\Delta h-(k+1)\lambda_1^2h^2\}\ge k\Delta h/2.
\]
Also $1-kd_i<1$. Appendix~\ref{app:privacy} gives each separate privacy
statement with exactly \eqref{eq:privacy-cal}. The prescribed guard is
$\epsilon(a_{\rm priv}-1)h/(4Hq^2)\le1$ for the source, and $1/k$ times
that quantity for the target. Finally, the explicit expressions in
\eqref{eq:constants} have the stated order when the public parameters are
fixed. The memory error is $O(\sqrt h)$.
\end{proof}

\paragraph{Limiting second memory.}
For every fixed $t>0$, \eqref{eq:second-L2} with
$l=\lfloor t/\tau\rfloor$ proves coordinatewise $L^2$ convergence of
$W_l=\tau\widehat v_l$ to $c^2$ as $\tau\to0$.
Thus the limiting denominator is the isotropic constant $c+e_0$.
The parameter $\lambda_2$ affects the finite error but disappears from
\eqref{eq:limit}. The result establishes a sufficient recipe in this
specified noise regime, not its necessity or a general active-clipping theorem.

\subsection{Sharper constants for the exponential memories}\label{sec:sharp}

The factors $E_i=e^{\lambda_iH}$ in \eqref{eq:constants} can be avoided.
Put $P_i=1+\lambda_iH/2$, $a=\sqrt{P_1}$, $b=\sqrt{P_2}$ and define
\begin{align}
 B_\tau^*&=b(\sqrt2c^2+2cM\sqrt\tau),\notag\\
 F_\tau^*&=\frac p{ce_0D}\left[
 2MB_\tau^*\sqrt{H\tau}
 +caB_\tau^*\sqrt\tau(1+\log n_\tau)
 +HM^3\tau+2caM^2\tau\sqrt H\right],\notag\\
 S_\tau^*&=\frac{2Ma\sqrt{H\tau}}D,\qquad
 A_\tau^*=e^{L_0H/D}(F_\tau^*+S_\tau^*),\notag\\
 Q_\tau^*&=2M\tau(1+\lambda_1H+\log n_\tau)
 +\sqrt p c\sqrt\tau(\lambda_1H+\log n_\tau+5a)
 +4\sqrt p c\lambda_1a\tau\sqrt H,\notag\\
 \omega_\tau^*&=\frac{M\tau+2\sqrt p ca\sqrt\tau}D,\qquad
 C_\tau^*=e^{L_0H/D}\left(\frac{Q_\tau^*}D+
                                     \frac{L_0H\omega_\tau^*}D\right),\notag\\
 \mathcal E^*(\tau;M,L_0,L_\phi)&=L_\phi(A_\tau^*+C_\tau^*+\omega_\tau^*).
 \label{eq:sharp-constants}
\end{align}

\begin{proposition}\label{prop:sharp}
Theorem~\ref{thm:coupling} remains true with the starred constants.
Consequently either occurrence of $\mathcal E$ in \eqref{eq:main-error}
may independently be replaced by $\mathcal E^*$, or by their minimum.
For the Poisson oracle on a reference dataset of $N$ records, one may further
replace $M$ in $S_\tau$ or $S_\tau^*$ only by
$\nu_q=\sqrt{(1-q)/(qN)}$.
\end{proposition}

\begin{proof}
The sharper weight, convolution and quadrature estimates are proved in
Appendix~\ref{app:sharp}. They replace the corresponding estimates in
Appendices~\ref{app:moments}--\ref{app:sde}; the dependence-preserving
Cauchy--Schwarz and Gronwall steps are unchanged. For the final assertion,
independence of the current record indicators, conditional on the past, gives
\[
 \E[\norm{A_j(x_j)-\nabla f(x_j)}^2\mid\mathcal F_j]
 =\frac{1-q}{qN^2}\sum_{r=1}^N\norm{g_r(x_j)}^2
 \le\frac{1-q}{qN}.
\]
Deterministic-weight orthogonality replaces the sampling estimate by
$\nu_q a/\sqrt l$ (or $\nu_q a_1/\sqrt l$), proving precisely the claimed
change in $S$. It does not change the signal bounds used elsewhere.
At $q=1$ this sampling contribution vanishes exactly.
\end{proof}

\section{A finite example and the epoch boundary}\label{sec:example}

Take four identical scalar records
$f_r(x)=\sqrt{1+x^2}-1$, $x_0=1$, and
\[
 p=H=e_0=\lambda_1=\lambda_2=L=1,\quad k=2,\quad
 q_s=\tfrac14,\quad q_t=\tfrac12,\quad
 \epsilon=\tfrac12,\quad\delta=e^{-4},\quad n=10^{20}.
\]
Here $a_{\rm priv}=17$, $c=\sqrt{17}/2$, $h=1/(2n)$, and
$\bar h=1/32$ satisfies all guards. The source and target privacy
small-step quantities are $32h$ and $16h$. Each record gradient has norm
at most one and Lipschitz constant one.

For either mesh, $\tau\le10^{-6}$. Use $2<c<21/10$, $D>3$,
$a=b=\sqrt{3/2}<5/4$, $\sqrt2<3/2$ and $M\le4$.
The exponential series gives $e^{1/D}<e^{1/3}<3/2$, by comparison
with the geometric series. Thus $B_\tau^*<25/3$.
Substitute these bounds and $\sqrt\tau\le1/1000$ into
\eqref{eq:sharp-constants}, retaining terms proportional to
$\tau$ with their extra factor $\sqrt\tau$. Direct rational arithmetic gives
\[
 \frac{\mathcal E^*(\tau;M,1,2)}{\sqrt\tau}
 <\frac{224551}{3000}+\frac{26091}{2000}\log n_\tau
 <80+14\log n_\tau.
\]
No improved Poisson sampling variance is needed for this estimate.
Since $e^3>10$ and $e>2$, $\log n<60$ and $\log(2n)<61$.
The sum of the two exponential-run errors is below $1854/10^{10}$.
For the target memory comparison, $K=1$, $C_\beta=3/2$, $C_x=4$,
and $C_\beta\bar h+HC_x<5$. Using $e<3$ gives the risk error
$729\sqrt{h/2}=729/(2\cdot10^{10})$. Consequently
\begin{equation}\label{eq:finite-example}
 |R_t-R_s|<\frac{4437}{2\cdot10^{10}}<10^{-6}.
\end{equation}
This extremely large finite example demonstrates nonempty quantitative
scope; it is not a computational recommendation or a bound forcing either
risk to be small.

\paragraph{Fixed epochs.}
The calibration used here requires
$\epsilon(a_{\rm priv}-1)/(4Tq^2)\le1$.
If the expected epoch count $E=Tq$ is fixed, this is equivalent to
\[
 q\ge\frac{\epsilon(a_{\rm priv}-1)}{4E},\qquad
 T\le\frac{4E^2}{\epsilon(a_{\rm priv}-1)}.
\]
It cannot justify this noise schedule along $q\to0$ at fixed $E$.
Moreover the error constants contain $M=1/q$, and the memory comparison
contains $L_q$. Uniform limits with bounded epochs require new calibration
and stability estimates. Neither the finite example nor the fixed-rate
asymptotic statement supplies them. Releasing the two runs jointly requires
privacy composition in addition to the separate conclusions above.

\appendix
\section{Same-release moments and the mean recursion}\label{app:moments}

At age $l\ge1$, define the deterministic corrected weights
\[
 w_{i,l,r}=\frac{e^{-\lambda_i\tau r}}
                  {\sum_{s=0}^{l-1}e^{-\lambda_i\tau s}},\quad 0\le r<l.
\]
They sum to one and satisfy
\begin{equation}\label{eq:weight-squares}
 \max_r w_{i,l,r}\le E_i/l,\qquad \sum_rw_{i,l,r}^2\le E_i/l,
 \qquad l\tau\le H.
\end{equation}
The ratio of largest to smallest unnormalized weight is at most $E_i$,
which proves both assertions. Let $W_l=\tau\widehat v_l$.
Multiplying the denominator in \eqref{eq:adam} by $\sqrt\tau$ gives exactly
\begin{equation}\label{eq:scaled-update}
 x_l=x_{l-1}-\tau\frac{\widehat m_l}{\sqrt{W_l}+e_0}.
\end{equation}
In particular its coefficient is $\tau$, including at $l=1$.

For a coordinate $i$, Minkowski's inequality gives
\begin{equation}\label{eq:first-L2}
 \Ltwo{\widehat m_{l,i}}\le M+\frac{ca_1}{\sqrt{l\tau}}.
\end{equation}
For its second memory, expand under the actual adapted law:
\begin{align*}
 W_{l,i}-c^2={}&c^2\sum_{j<l}w_{2,l,l-1-j}(Z_{j,i}^2-1)\\
 &+2c\sqrt\tau\sum_{j<l}w_{2,l,l-1-j}A_{j,i}(x_j)Z_{j,i}
 +\tau\sum_{j<l}w_{2,l,l-1-j}A_{j,i}(x_j)^2.
\end{align*}
The first sum has variance $2\sum w_2^2$. The middle summands are
orthogonal martingale differences: their signal is known before the current
Gaussian and may depend on earlier releases. The last sum is bounded by
$\tau M^2$. Minkowski and \eqref{eq:weight-squares} therefore give
\begin{equation}\label{eq:second-L2}
 \Ltwo{W_{l,i}-c^2}\le B_\tau/\sqrt l+\tau M^2.
\end{equation}
No independence between the three sums is used.

For every $W\ge0$,
\[
 \left|\frac1{\sqrt W+e_0}-\frac1D\right|
 =\frac{|W-c^2|}{(\sqrt W+c)(\sqrt W+e_0)D}
 \le\frac{|W-c^2|}{ce_0D}.
\]
Let $r_l=\widehat m_l/(\sqrt{W_l}+e_0)-\widehat m_l/D$.
Cauchy--Schwarz on each coordinate, followed by the sum of coordinate
absolute values, yields
\begin{equation}\label{eq:den-error}
 \E\norm{r_l}\le\frac p{ce_0D}
 \left(M+\frac{ca_1}{\sqrt{l\tau}}\right)
 \left(\frac{B_\tau}{\sqrt l}+\tau M^2\right).
\end{equation}
The first and second memories remain dependent throughout this calculation.
Using $\sum_{l=1}^n l^{-1/2}\le2\sqrt n$ and
$\sum_{l=1}^n l^{-1}\le1+\log n$ proves
\begin{equation}\label{eq:den-sum}
 \tau\sum_{l=1}^{n_\tau}\E\norm{r_l}\le F_\tau.
\end{equation}

Define an auxiliary sequence $z$ with the same Gaussian innovations,
corrected first-memory weights and zero initialization, using releases
$b(z_j)+cZ_j/\sqrt\tau$ and update
$z_l=z_{l-1}-\tau\widehat m_l^z/D$.
For the actual sequence, set $\varepsilon_j=A_j(x_j)-b(x_j)$.
Conditional variance gives
$\E[\norm{\varepsilon_j}^2\mid\mathcal F_j]\le M^2$.
These are martingale differences at the actual positions, and deterministic
weighted orthogonality implies
\[
 \E\norm{\sum_{j<l}w_{1,l,l-1-j}\varepsilon_j}\le Ma_1/\sqrt l.
\]
In the difference of the two first memories the Gaussian terms cancel.
The remaining drift difference is at most
$L_0\max_{j<l}\norm{x_j-z_j}$. Hence for
$\mathfrak e_l=\max_{j\le l}\norm{x_j-z_j}$,
\[
 \mathfrak e_l\le(1+\tau L_0/D)\mathfrak e_{l-1}+\tau\norm{r_l}
    +\frac\tau D\norm{\sum_{j<l}w_{1,l,l-1-j}\varepsilon_j}.
\]
Iteration, expectations and \eqref{eq:den-sum} prove
$\E \mathfrak e_{n_\tau}\le A_\tau$. This stability estimate uses the mean
Lipschitz constant, rather than the worst sampled Lipschitz constant.

\section{Brownian comparison including the first interval}\label{app:sde}

Write $\lambda=\lambda_1$, $E=E_1$ and couple
$Z_j=(B_{(j+1)\tau}-B_{j\tau})/\sqrt\tau$. Independent Brownian bridges
give such a coupling without changing the release law. For a bounded
adapted input $a$ with $\norm a\le M$, put
\begin{align*}
 r_a(t)&=\frac\lambda{1-e^{-\lambda t}}
              \int_0^te^{-\lambda(t-s)}a(s)\,ds,\\
 \xi(t)&=\frac{\lambda c}{1-e^{-\lambda t}}
              \int_0^te^{-\lambda(t-s)}\,dB_s.
\end{align*}
The first kernel has mass one. Since
$1-e^{-\lambda t}\ge\lambda t/E$, It\^o isometry gives
\begin{equation}\label{eq:noise-start}
 \Ltwo{\xi(t)}\le\sqrt p cE/\sqrt t.
\end{equation}
In particular $\int_0^H\norm{\xi(t)}dt$ is integrable and finite almost
surely, and its indefinite integral is continuous at zero.

Take $a(s)=b(z_j)$ on $[j\tau,(j+1)\tau)$. At $t_l=l\tau$,
$r_a(t_l)$ is exactly the corrected deterministic part of
$\widehat m_l^z$. The unnormalized discrete Gaussian kernel on each
interval is the interval average of
$\lambda e^{-\lambda(t_l-s)}$. The derivative of this function in $s$
is at most $\lambda^2$, so the pointwise difference from its average is
at most $\lambda^2\tau$. After normalization and It\^o isometry,
\begin{equation}\label{eq:gaussian-kernel-error}
 \E\norm{\widehat m_l^z-r_a(t_l)-\xi(t_l)}
 \le\sqrt p c\lambda E\tau/\sqrt{t_l}.
\end{equation}
The sum of these bounds times $\tau$ is at most
$2\sqrt p c\lambda E\tau\sqrt H$.
For clarity, the two exact cell identities behind this calculation are
\[
 \int_{j\tau}^{(j+1)\tau}\lambda e^{-\lambda(t_l-s)}ds
  =(1-e^{-\lambda\tau})e^{-\lambda\tau(l-1-j)},\qquad
 \frac{w_{1,l,l-1-j}}{\tau}
  =\frac{1}{\tau(1-e^{-\lambda t_l})}
       \int_{j\tau}^{(j+1)\tau}\lambda e^{-\lambda(t_l-s)}ds.
\]
Thus the Gaussian kernel is a cell average, while the drift cell integral
equals the corrected weight. Neither identity discards the finite
bias-correction factor.

For $t>0$, the locally absolutely continuous deterministic convolution obeys
$r_a'(t)=\lambda(a(t)-r_a(t))/(1-e^{-\lambda t})$ almost everywhere.
Its derivative norm is at most $2ME/t$; it is also bounded by $M$.
Its integrated right-endpoint quadrature error is therefore at most
$2M\tau$ on the first interval and $2ME\tau\log n_\tau$ on the rest.

For $0<s\le t\le H$, split the Gaussian integral at $s$. The multiplier
of the old noise relative to $\xi(s)$ is
\[
 r=\frac{e^{\lambda s}-1}{e^{\lambda t}-1},\qquad
 0\le1-r\le E(t-s)/s.
\]
The fresh integral and \eqref{eq:noise-start} give
\begin{equation}\label{eq:noise-modulus}
 \E\norm{\xi(t)-\xi(s)}\le\sqrt p c
 \left(\frac{E\sqrt{t-s}}s+\frac{E^2(t-s)}{s^{3/2}}\right).
\end{equation}
On the first interval \eqref{eq:noise-start} bounds the integrated
right-endpoint error by $3\sqrt p cE\sqrt\tau$. Elsewhere set $t$ to
the right endpoint, use $t-s\le\tau$, and integrate
\eqref{eq:noise-modulus}. The resulting bound is
$\sqrt p c\sqrt\tau\{E(3+\log n_\tau)+2E^2\}$.
Together with \eqref{eq:gaussian-kernel-error} this proves
\begin{equation}\label{eq:integral-residual}
 \E\max_{l\le n_\tau}\left\|z_l-x_0+
       D^{-1}\int_0^{l\tau}(r_a(s)+\xi(s))ds\right\|\le Q_\tau/D.
\end{equation}
Each endpoint discrepancy is a partial sum of interval errors. Bounding
its maximum by the sum of absolute interval errors justifies the order
of the maximum and expectation in this formula.

For each path with integrable $\xi$, solve
$X_t=x_0-D^{-1}\int_0^t(r_{b(X)}(s)+\xi(s))ds$ by Picard iteration.
The normalized drift kernel satisfies
\[
 \norm{r_{b(X)}(s)-r_{b(Y)}(s)}
       \le L_0\sup_{u\le s}\norm{X_u-Y_u}.
\]
Successive differences are bounded by the factorial series for
$\exp(L_0H/D)$, since the first difference is bounded on $[0,H]$.
This gives a continuous solution; Gronwall gives uniqueness. The iterations
are causal, hence the solution is adapted. Defining $m$ by
\eqref{eq:limit} proves that equation with zero initialization.
In addition, the integral equation gives the useful domination
\[
 \E\sup_{0\le t\le H}\norm{X_t-x_0}
 \le \frac{MH+2\sqrt p\,cE\sqrt H}{D}<\infty.
\]
This justifies the subsequent expectations and time-integral interchanges,
together with the preceding nonnegative integral bounds.
The bound $\norm{r_{b(X)}}\le M$ and \eqref{eq:noise-start} imply
\begin{equation}\label{eq:X-modulus}
 \E\norm{X_t-X_s}\le\omega_\tau\quad(0\le s\le t\le H,\ t-s\le\tau),
\end{equation}
using $\int_s^t u^{-1/2}du\le2\sqrt{t-s}$.

Set $\mathfrak e_l=\max_{j\le l}\norm{z_j-X_{j\tau}}$. At time
$u\in[j\tau,(j+1)\tau)$,
$\norm{b(z_j)-b(X_u)}\le L_0(\mathfrak e_j+\norm{X_{j\tau}-X_u})$.
Every earlier grid index is bounded by $\mathfrak e_j$. Positivity and mass one of
the normalized kernel, \eqref{eq:X-modulus}, and
\eqref{eq:integral-residual} thus give
\[
 \E \mathfrak e_l\le Q_\tau/D+L_0H\omega_\tau/D
               +(\tau L_0/D)\sum_{j<l}\E \mathfrak e_j.
\]
Discrete Gronwall proves $\E \mathfrak e_{n_\tau}\le C_\tau$, completing the
coupling argument. The within-interval bound was used inside expectation;
no pathwise Lipschitz bound on Brownian paths was assumed.

\section{Perturbing both corrected memory parameters}\label{app:memory}

For $b\in(0,1)$ let $w_r(b)=b^r/\sum_{s<l}b^s$.
If $a,b\in(0,1)$ and $R=(l-1)|\log a-\log b|$, then
\begin{equation}\label{eq:relative-weights}
 e^{-R}\le w_r(a)/w_r(b)\le e^R.
\end{equation}
Indeed the ratio of the largest and smallest unnormalized ratios $(a/b)^r$
is $e^R$, and their $w(b)$-weighted average is the ratio of normalizers.
For exponential parameters at a mesh $\tau$ with $l\tau\le H$,
\[
 \sum_r\frac{w_r(e^{-\lambda_1\tau})^2}
                {w_r(e^{-\lambda_2\tau})}\le K^2.
\]
Weighted Cauchy--Schwarz therefore gives $|\widehat m_i|\le K\sqrt{\widehat v_i}$
for every scalar release sequence, including zero second memory.

Let $D_b(y)$ denote the scalar direction at a positive stabilizer $\zeta$.
When both parameters are changed to $a_i$, put
$R_i=(l-1)|\log a_i-\log b_i|$. Subtracting the two ratios, first
changing the numerator, and applying \eqref{eq:relative-weights}, gives
\begin{equation}\label{eq:beta-direction}
 |D_a(y)-D_b(y)|\le K(e^{R_1+R_2/2}-1).
\end{equation}
In detail, the numerator difference is at most
$K(e^{R_1}-1)\sqrt{\widehat v_b}$; the root-memory ratio lies in
$[e^{-R_2/2},e^{R_2/2}]$. The two ratio differences are bounded by
$Ke^{R_2/2}(e^{R_1}-1)$ and $K(e^{R_2/2}-1)$. If the original root
memory vanishes, all releases vanish and the claim holds directly.

At fixed exponential parameters, for two vector histories with
$v=\max_{j<l}\norm{y_j-y'_j}$,
\begin{equation}\label{eq:history-direction}
 \norm{D_b(y)-D_b(y')}\le(1+K)v/\zeta.
\end{equation}
The first-memory difference is at most $v$. The reverse triangle
inequality in each weighted Euclidean space bounds the vector of root-memory
differences by $v$. In the ratio decomposition, each original normalized
numerator is at most $K$ and each denominator at least $\zeta$.
This proves \eqref{eq:history-direction} without independence assumptions.

For the target pair put $a_i=1-k(1-e^{-\lambda_i h})$ and
$b_i=e^{-\lambda_i kh}$. Taylor's theorem for $u\mapsto u^k$ on $[0,1]$
gives
\[
 0\le b_i-a_i\le\tfrac12k(k-1)\lambda_i^2h^2,
 \qquad 0\le\log b_i-\log a_i\le k(k-1)\lambda_i^2h^2,
\]
using $a_i\ge1/2$. For $l\le n$, this implies
$R_1+R_2/2\le C_\beta h$.
Couple the two target runs by the same masks and Gaussians. Their releases
differ by at most $L_q$ times their position difference; they need not be
identical. Apply \eqref{eq:beta-direction} to the linear run's history,
then \eqref{eq:history-direction} to change that history. For the running
position maximum $\mathfrak e_j$, the target step and stabilizer give
\[
 \mathfrak e_{j+1}\le(1+khC_x)\mathfrak e_j
       +\sqrt{kh}\,\sqrt p K C_\beta e^{C_\beta\bar h}h.
\]
Since $nkh=H$, iteration yields
$\mathfrak e_n\le A_{\rm mem}\sqrt{h/k}$, pathwise. At the first update the
normalized weights coincide exactly. At $k=1$ the entire error is zero.

\section{Finite privacy for the adaptive transcripts}\label{app:privacy}

We prove the particular sufficient calibration used above, for public
algorithm parameters and normalizers fixed on all neighboring datasets.
All record gradients are Borel and bounded by one. Every finite memory and
position is a Borel function of preceding releases, because the stabilizer
is positive and both correction factors are nonzero. The release kernels
therefore have positive Borel densities.

For an adjacent pair $E,E\cup\{r\}$, condition on an arbitrary common
release history and the sampled mask of the common records. The state
$x$ is the same deterministic function of that history on both sides.
Let $b$ denote the density of the noisy common sum, and let $p$ denote its
translate by $g_r(x)$. The two conditional densities are $b$ and
$m=(1-q)b+qp$. The differing-record inclusion bit is not revealed.
The common mask can be included as an auxiliary variable: its law is the
same on both sides. Integrating it out and dividing sums by the fixed $B$
gives the actual release law at each age, hence the entire law by induction.

For integer $a\ge2$, Gaussian square completion gives
\begin{equation}\label{eq:mixture-power}
 \int m^a b^{1-a}=A_a(q,z):=
 \sum_{j=0}^a\binom aj(1-q)^{a-j}q^j e^{j(j-1)z/2},
 \qquad z=\norm{g_r(x)}^2/\sigma^2.
\end{equation}
The reverse power is no larger. Here is a proof including $q=1$.
For $u>0$ put $y=1-q+qu$, $w=1-q+q/u$ and
$F_a(u)=y^a+uw^a-y^{1-a}-uw^{1-a}$.
For $u\ge1$, $yw\ge1$, $y\ge1\ge w$, $y-1=u(1-w)$,
$F_1=0$, and
\[
 F_{a+1}-F_a=(y-1)\{y^a+y^{-a}-w^a-w^{-a}\}\ge0.
\]
The last inequality follows from $y\ge1/w$ and monotonicity of
$s^a+s^{-a}$ on $s\ge1$. For $u<1$, use $F_a(u)=uF_a(1/u)$.
Reflection about the midpoint of the two Gaussian means swaps $b,p$.
Pairing reflected points with $u=p/b$ now gives
$\int b^am^{1-a}\le\int m^ab^{1-a}$. All terms in
\eqref{eq:mixture-power} are nondecreasing in $z$, so $z\le\sigma^{-2}$
can be substituted.

If $a(a-1)/(2\sigma^2)\le1$, use $e^v\le1+2v$ for $v\in[0,1]$.
For $J\sim\operatorname{Binomial}(a,q)$ this gives
\[
 A_a(q,\sigma^{-2})
 \le1+\E[J(J-1)]/\sigma^2
 =1+a(a-1)q^2/\sigma^2.
\]
Thus each ordered row R\'enyi divergence is at most $aq^2/\sigma^2$.
The bound holds at every history and can be averaged over the common mask.
The marginalization step is an inequality, not an identification of the
mixture divergence with the average conditional divergence. Explicitly,
let a finite mask $z$ have common law $\pi_z$, and let $p_z,q_z$ be the
positive conditional densities in either ordered neighboring direction.
Set $\bar p=\sum_z\pi_zp_z$ and $\bar q=\sum_z\pi_zq_z$. At every output,
Jensen's inequality for $u\mapsto u^a$ gives
\begin{align*}
 \bar p^{\,a}\bar q^{\,1-a}
 &=\bar q\left(\sum_z\frac{\pi_zq_z}{\bar q}
                         \frac{p_z}{q_z}\right)^a\\
 &\le\sum_z\pi_zp_z^{\,a}q_z^{\,1-a}.
\end{align*}
Integration preserves the uniform conditional moment bound. Division by
the same positive public $B$ preserves the power integral by change of
variables. The differing-record inclusion bit is still unobserved; it was
integrated into $m=(1-q)b+qp$ before this common-mask calculation.
For two prefix densities $P_j,Q_j$ and conditional row densities $p_j,q_j$,
Tonelli gives
\[
 \int(P_jp_j)^a(Q_jq_j)^{1-a}
 \le e^{(a-1)\rho_j}\int P_j^aQ_j^{1-a},\qquad
 \rho_j=aq^2/\sigma^2.
\]
Induction proves the ordered transcript bound with the sum of $\rho_j$;
the same proof with the laws reversed covers the other orientation.
Under \eqref{eq:privacy-cal} this sum is $\epsilon/2$.
For the transcript log likelihood $L=\log(P/Q)$, Markov's inequality gives
\[
 P\left\{L>\epsilon/2+\frac{\log(1/\delta)}{a_{\rm priv}-1}\right\}
 \le\delta.
\]
The threshold is at most $\epsilon$. Splitting any measurable event on
this exceptional set proves $P(S)\le e^\epsilon Q(S)+\delta$;
reverse the laws for the opposite inequality. The same deterministic
projection or independent public randomization preserves these inequalities.
This proves privacy for the actual initialized trajectories and selected
parameters, separately for each run, with no independence assumption on
successive adaptive releases.

\section{Exact normalized-memory estimates}\label{app:sharp}

For $x>0$, $\tanh x\le x$ and
$\coth x=1+2/(e^{2x}-1)\le1+1/x$. The corrected exponential weights
therefore obey the exact identity and bound
\begin{equation}\label{eq:sharp-weights}
 \sum_{r<l}w_{i,l,r}^2
 =\tanh(\lambda_i\tau/2)\coth(\lambda_i l\tau/2)
 \le\frac1l+\frac{\lambda_i\tau}2\le\frac{P_i}l.
\end{equation}
This bounds their squared sum; no corresponding replacement of the
maximum-weight bound is needed. It proves the starred versions of
\eqref{eq:first-L2}--\eqref{eq:den-sum} and the sampling estimate.

Write $\lambda=\lambda_1$. Direct It\^o isometry gives
\begin{equation}\label{eq:sharp-noise}
 \E\norm{\xi(t)}^2
 =pc^2\frac\lambda2\coth(\lambda t/2)
 \le pc^2(1/t+\lambda/2)\le pc^2P_1/t.
\end{equation}
In particular the first interval and the solution's time modulus are
controlled as before, now with $a$ in place of $E_1$.

For a continuously differentiable function $g$ on an interval $I$ of length
$\tau$, let $\bar g$ be its average. The variance identity and
Cauchy--Schwarz show
\[
 \int_I|g-\bar g|^2
 =\frac1{2\tau}\int_I\int_I|g(s)-g(r)|^2\,ds\,dr
 \le\tau^2\int_I|g'|^2.
\]
At a grid time the discrete Gaussian kernel is the cellwise average of
the continuous kernel $g(s)=\lambda e^{-\lambda(t_l-s)}/(1-e^{-\lambda t_l})$.
Since $g'=\lambda g$, summing the displayed estimate over all cells
and using \eqref{eq:sharp-noise} gives
\[
 \E\norm{\widehat m_l^z-r_a(t_l)-\xi(t_l)}
 \le\sqrt p c\lambda a\tau/\sqrt{t_l}.
\]
The sum times $\tau$ is at most $2\sqrt p c\lambda a\tau\sqrt H$.

The elementary bound
$\lambda/(1-e^{-\lambda t})\le\lambda+1/t$
implies $\norm{r_a'(t)}\le2M(\lambda+1/t)$.
Its first-interval quadrature error is at most $2M\tau$; all later
intervals together cost at most $2M\tau(\lambda H+\log n_\tau)$.

For the Gaussian quadrature let $0<s\le t\le H$ and $\Delta=t-s$.
The old-noise multiplier $r$ from Appendix~\ref{app:sde} satisfies
\[
 1-r=\frac{1-e^{-\lambda\Delta}}{1-e^{-\lambda t}}
 \le\Delta(\lambda+1/t)\le\Delta(\lambda+1/s).
\]
The fresh Gaussian integral has $L^2$ norm at most
$\sqrt p c\sqrt\Delta(\lambda+1/s)$. Hence
\[
 \E\norm{\xi(t)-\xi(s)}
 \le\sqrt p c\left[
 \sqrt\Delta(\lambda+1/s)
 +a\Delta\left(\frac\lambda{\sqrt s}+\frac1{s^{3/2}}\right)\right].
\]
On the first interval the integrated right-endpoint error is at most
$3\sqrt p ca\sqrt\tau$. Integrating the last display over
$s\in[\tau,H]$, with $t$ the right mesh endpoint and $\Delta\le\tau$,
gives at most
\[
 \sqrt p c\sqrt\tau(\lambda H+\log n_\tau+2a)
 +2\sqrt p c\lambda a\tau\sqrt H.
\]
Adding the first interval, the Gaussian-kernel comparison, and the
deterministic quadrature proves $Q_\tau^*$ in \eqref{eq:sharp-constants}.
The same summation of absolute interval errors controls the maximum over
all grid endpoints. The remaining Volterra and observable comparisons
are unchanged. When $n_\tau=1$, the later-interval integrals are empty;
the displayed upper bounds remain valid.

\section{A sharper expected memory-transfer estimate}\label{app:expected-memory}

The pathwise estimate in Appendix~\ref{app:memory} applies to arbitrary
release histories. For the actual bounded-signal Gaussian releases, keeping
the stabilizer in the estimate gives an additional expected bound. This
does not alter Theorem~\ref{thm:pair}, its hypotheses, or its finite witness.

\begin{proposition}[Expected memory transfer]\label{prop:expected-memory}
Under the assumptions of Theorem~\ref{thm:pair}, let $x^{\rm lin}$ and
$x^{\rm exp}$ be the coupled linear and exponential target runs. Put
$M_t=1/(kq)$ and $a_1=e^{\lambda_1H/2}$, and define
\begin{equation}\label{eq:expected-memory-constant}
 B_{\rm mem}=\frac{C_\beta}{e_0}
       e^{C_\beta\bar h+HC_x}
       \left(M_tH+2\sqrt p\,c a_1\sqrt H\right).
\end{equation}
Then
\begin{equation}\label{eq:expected-memory}
 \E\max_{0\le j\le n}\norm{x_j^{\rm lin}-x_j^{\rm exp}}
 \le B_{\rm mem}h.
\end{equation}
Consequently the right side of \eqref{eq:main-error} may be strengthened to
\begin{equation}\label{eq:sharper-pair}
 2L\min\{A_{\rm mem}\sqrt{h/k},B_{\rm mem}h\}
 +\mathcal E(h;1/q,L,2L)+\mathcal E(kh;1/(kq),L,2L).
\end{equation}
\end{proposition}

\begin{proof}
Write $\tau=kh$, $\zeta=e_0/\sqrt\tau$,
$\alpha_i=1-k(1-e^{-\lambda_i h})$, and
$\gamma_i=e^{-\lambda_i\tau}$. At age $l\ge1$, set
$R_i=(l-1)|\log\alpha_i-\log\gamma_i|$ and
\[
 U_l=M_t+\frac{\sqrt p\,c a_1}{\sqrt{l\tau}}.
\]
All the weights below are deterministic. Work first on the one actual
linear-run release history $Y_j=A_j+cZ_j/\sqrt\tau$, where
$\norm{A_j}\le M_t$. The signals may depend on earlier Gaussian
innovations. Denote by $m_\alpha,m_\gamma$ the corrected first moments
computed with the two parameters on this same history. The exponential
weights satisfy $\sum w(\gamma_1)^2\le a_1^2/l$, so Minkowski and the
Gaussian second moment give
\[
 \E\norm{m_\gamma}\le U_l.
\]
For $d_j=w_{l,l-1-j}(\alpha_1)-w_{l,l-1-j}(\gamma_1)$,
\eqref{eq:relative-weights} gives
$|d_j|\le(e^{R_1}-1)w_{l,l-1-j}(\gamma_1)$. Therefore
\begin{align*}
 \E\norm{m_\alpha-m_\gamma}
 &\le M_t\sum_j|d_j|
    +\frac c{\sqrt\tau}\left(\E\norm{\sum_jd_jZ_j}^2\right)^{1/2}\\
 &=M_t\sum_j|d_j|+\frac{\sqrt p\,c}{\sqrt\tau}
                      \left(\sum_jd_j^2\right)^{1/2}
 \le(e^{R_1}-1)U_l.
\end{align*}
This bounds the signed Gaussian weight difference directly. It does not
assume independence between the adapted signal sum and the Gaussian sum.

Let $r_\alpha,r_\gamma$ denote the coordinatewise root second moments.
Their relative weight bounds imply, in each coordinate,
\[
 \frac{|r_\alpha-r_\gamma|}{r_\alpha+\zeta}
 \le e^{R_2/2}-1.
\]
When $r_\gamma=0$, both roots and numerators vanish; otherwise use
$e^{-R_2/2}r_\gamma\le r_\alpha\le e^{R_2/2}r_\gamma$ and split
according to which root is larger. Decomposing the two directions as
in Appendix~\ref{app:memory}, but retaining the lower denominator bound
$\zeta$, gives the pointwise vector inequality
\[
 \norm{D_\alpha(Y)-D_\gamma(Y)}
 \le\frac{\norm{m_\alpha-m_\gamma}
             +(e^{R_2/2}-1)\norm{m_\gamma}}{\zeta}.
\]
Taking expectations and using
$(e^u-1)+(e^v-1)\le e^{u+v}-1$ for $u,v\ge0$ yields
\[
 \E\norm{D_\alpha(Y)-D_\gamma(Y)}
 \le\frac{e^{R_1+R_2/2}-1}{\zeta}U_l
 \le\frac{C_\beta e^{C_\beta\bar h}h}{\zeta}U_l.
\]
There is no factorization of dependent random numerator and denominator
terms.

Let $\mathfrak e_l=\max_{j\le l}\norm{x_j^{\rm lin}-x_j^{\rm exp}}$.
The pathwise history comparison \eqref{eq:history-direction} still gives
\[
 \mathfrak e_l\le(1+\tau C_x)\mathfrak e_{l-1}
       +\sqrt\tau\,\norm{D_\alpha(Y^{\rm lin})
                                -D_\gamma(Y^{\rm lin})}.
\]
Iterate from $\mathfrak e_0=0$, take expectations, and use
$\sqrt\tau/\zeta=\tau/e_0$. Since
\[
 \sum_{l=1}^n\tau U_l
 \le M_tH+2\sqrt p\,c a_1\sqrt H,
\]
the result is \eqref{eq:expected-memory}. The Lipschitz observable and the
two already-proved memory bounds give \eqref{eq:sharper-pair}.
\end{proof}

\begin{remark}[Consequences and limitations]
At $k=1$, $C_\beta=B_{\rm mem}=0$. At $n=1$ the corrected weights agree
and both target first updates coincide. Legal full target sampling $kq=1$
also satisfies the proof. For the witness in Section~\ref{sec:example},
$M_t=2$, $c<3$, $a_1<2$ and
$e^{C_\beta\bar h+HC_x}<3^5$ imply $B_{\rm mem}<5103$.
At the current $n=10^{20}$ example its memory contribution is therefore
below $5103/10^{20}$. Combining this with the sharper single-run bounds
in Section~\ref{sec:sharp} gives the complete comparison
\[
 |R_t-R_s|<\frac{1854}{10^{10}}+\frac{5103}{10^{20}}<2\cdot10^{-7}.
\]
The earlier $n=10^{24}$ example with the original coupling constants
instead gives $294000/10^{12}+5103/10^{24}$, also below $10^{-6}$.
The complete rate remains $O(\sqrt h(1+\log(1/h)))$.
The squared-weight estimate \eqref{eq:sharp-weights} also permits replacing
$a_1$ in $B_{\rm mem}$ by $\sqrt{1+\lambda_1H/2}$: the proof uses
$a_1$ only in the two deterministic-weight Gaussian second-moment estimates.
All other factors remain unchanged.
The factors $M_t$ and $C_x$ still depend on $1/q$, and the privacy guard
is unchanged, so this strengthening supplies no bounded-epoch uniformity.
\end{remark}

\section{A constructive comparison of privacy laws}\label{app:canonical}

This appendix supplies a reusable law comparison for the same adaptive
release mechanism. It does not replace \eqref{eq:privacy-cal}, and no
alternative calibration is applied to Theorem~\ref{thm:pair}.

For probability measures $P,Q$, define the hockey-stick divergence
\[
 H_\epsilon(P\Vert Q)=\sup_{A\text{ measurable}}
                       \{P(A)-e^\epsilon Q(A)\}.
\]

\begin{proposition}[One common kernel for an adaptive transcript]
\label{prop:canonical-kernel}
Fix an integer $T\ge1$, $0<q\le1$, $\sigma\ge0$, a sensitivity bound
$\Delta>0$, and a positive public normalizer $B$. Suppose each adaptive
state is the same Borel function of preceding releases and fixed public
parameters on every dataset,
and every record contribution is Borel and globally bounded in norm by
$\Delta$. Define the scalar product laws
\[
 P_{\rm can}=\bigl[(1-q)\mathcal N(0,\sigma^2)
                 +q\mathcal N(\Delta,\sigma^2)\bigr]^{\otimes T},
 \qquad Q_{\rm can}=\mathcal N(0,\sigma^2)^{\otimes T},
\]
where a zero-variance normal law is a point mass. For each adjacent pair
$(D_+=D_-\cup\{r\},D_-)$ there is one Borel Markov kernel $\mathsf K$
such that their actual Poisson-sampled Gaussian transcript laws satisfy
\[
 \mathcal L_{D_+}=P_{\rm can}\mathsf K,
 \qquad \mathcal L_{D_-}=Q_{\rm can}\mathsf K.
\]
Consequently each ordered actual hockey-stick divergence is no larger than
the corresponding ordered canonical-product divergence. A two-sided
$(\epsilon,\delta)$ guarantee for the canonical pair is therefore
sufficient for the actual adaptive transcript.
\end{proposition}

\begin{proof}
For a prefix $z$, let $x(z)$ be the state and put
$u(z)=g_r(x(z))/\Delta$, so $\norm{u(z)}\le1$. Define on the closed unit
ball
\[
 S(u)=I-\frac{uu^\top}{1+\sqrt{1-\norm u^2}}.
\]
This matrix-valued map is continuous, including at $u=0$ and $\norm u=1$.
Its denominator is at least one. It is symmetric and satisfies
\[
 S(u)^2=I-uu^\top.
\]
Indeed, with $s=\sqrt{1-\norm u^2}$, the coefficient of $uu^\top$ in
$S(u)^2-I$ is
\[
 -\frac{2}{1+s}+\frac{1-s^2}{(1+s)^2}=-1.
\]
Given a scalar input $U$ and independent $G\sim\mathcal N(0,I_p)$, set
\[
 V=u(z)U+\sigma S(u(z))G.
\]
If $U\sim\mathcal N(0,\sigma^2)$, the output has mean zero and covariance
$\sigma^2\{uu^\top+S(u)S(u)^\top\}=\sigma^2I_p$. If instead
$U\sim\mathcal N(\Delta,\sigma^2)$, its mean is $g_r(x(z))$ with the
same covariance. These are Gaussian laws, by closure under independent
linear combinations. The same channel consequently maps the scalar
mixture to the vector mixture with the full record shift and one hidden
inclusion bit.

Independently sample Bernoulli$(q)$ indicators $J_s$ for the common records
$s\in D_-$, form $C(z,J)=\sum_{s\in D_-}J_sg_s(x(z))$, and return
$(C(z,J)+V)/B$. Under the centered canonical input this is the smaller
dataset's conditional row law. Under the mixture input it is the larger
dataset's conditional row law. The masks, fresh Gaussian, and all auxiliary
randomness stay internal to the channel.

Apply this construction recursively to the $T$ scalar input coordinates,
using fresh auxiliary randomness at each age and computing the next state
from the generated prefix. All maps are Borel and the recursion is finite,
so it defines a Borel Markov kernel $\mathsf K$. Under either canonical
product law, the next input coordinate is independent of the previously
generated prefix. Induction therefore identifies the two output laws with
the actual adaptive transcripts. The kernel can depend on the fixed
neighboring pair, but the same kernel is used for both laws of that pair.
This is a comparison of probability laws, not a change to the released
mechanism.

For an output event $A$, let $f(z)=\mathsf K(z,A)\in[0,1]$.
Layer-cake integration gives, for any input laws $P,Q$,
\begin{align*}
 (P\mathsf K)(A)-e^\epsilon(Q\mathsf K)(A)
 &=\int_0^1\bigl(P\{f>t\}-e^\epsilon Q\{f>t\}\bigr)dt\\
 &\le H_\epsilon(P\Vert Q).
\end{align*}
Taking the supremum over $A$ and then reversing $P,Q$ proves both ordered
contractions.
\end{proof}

\begin{remark}[Scope of the comparison]
The construction preserves a common record mask across coordinates,
adaptive states, public normalization, and the hidden differing-record
bit. It includes $q=1$, zero shift, and maximal shift. A canonical-product
calibration is sufficient; it need not be the exact privacy requirement
of every adaptive transcript or of a selected parameter alone.
The proposition does not prove a new finite calibration, a sparse-rate
asymptotic, or a bounded-epoch utility comparison. Those require their own
estimates. For the current theorem, Appendix~\ref{app:privacy} remains the
complete proof of the prescribed guarded calibration.
\end{remark}

\end{document}
